
\documentclass[12pt,a4paper]{article}


\usepackage[utf8]{inputenc}
\usepackage[T1]{fontenc}
\usepackage[brazil]{babel}
\usepackage{geometry}
\usepackage{array}
\usepackage{longtable}
\usepackage{enumitem}
\usepackage{amsmath}
\usepackage{hyperref}


\geometry{
    a4paper,
    left=3cm,
    right=2cm,
    top=3cm,
    bottom=2cm
}


\hypersetup{
    colorlinks=true,
    linkcolor=black,
    urlcolor=blue
}

\begin{document}

\begin{center}
    {\LARGE \textbf{Sistema de Gestão de Biblioteca}}\\[0.3cm]
    {\Large Acervo Manager}
\end{center}

\vspace{1cm}

\section{Arquitetura}

A arquitetura do Acervo Manager define como o software será estruturado internamente, quais tecnologias serão utilizadas e como os diferentes componentes do sistema irão se comunicar.

A arquitetura foi planejada com uma separação entre frontend, backend e banco de dados, permitindo que cada componente tenha uma responsabilidade específica.

\section{Tecnologias utilizadas}

Para o desenvolvimento do sistema, foram definidas as seguintes tecnologias:

\begin{table}[h!]
    \centering
    \renewcommand{\arraystretch}{1.4}
    \begin{tabular}{|p{3.5cm}|p{4cm}|p{6cm}|}
        \hline
        \textbf{Componente} & \textbf{Tecnologia} & \textbf{Função} \\
        \hline
        Front-End & Angular & Desenvolvimento da interface do sistema \\
        \hline
        Estilização & Tailwind CSS & Criação e organização das telas do sistema \\
        \hline
        Back-End & Node.js + Express & Processamento das requisições e lógica do sistema \\
        \hline
        Banco de dados & MySQL & Armazenamento das informações \\
        \hline
        Comunicação & API REST / HTTP & Comunicação entre front-end e back-end \\
        \hline
    \end{tabular}
    \caption{Tecnologias utilizadas no Acervo Manager}
    \label{tab:tecnologias}
\end{table}

\section{Frontend}

O frontend será desenvolvido utilizando Angular, sendo responsável pela interface com a qual o usuário irá interagir.

O Angular será utilizado para desenvolver as diferentes telas e funcionalidades do sistema, como:

\begin{itemize}
    \item Cadastro e consulta de livros;
    \item Cadastro e consulta de usuários;
    \item Registro de empréstimos;
    \item Registro de devoluções;
    \item Consulta da disponibilidade dos livros.
\end{itemize}

Para a estilização da aplicação será utilizado o Tailwind CSS, permitindo criar uma interface organizada, responsiva e consistente.

\section{Backend}

O backend será desenvolvido utilizando Node.js com o framework Express.

Sua principal função será receber e processar as requisições realizadas pelo frontend, executar as regras necessárias do sistema e realizar a comunicação com o banco de dados.

Por exemplo, quando um usuário solicitar a lista de livros disponíveis, o Angular enviará uma requisição para a API. O backend receberá essa solicitação, consultará o MySQL e retornará os dados para o frontend.

\section{Comunicação entre Frontend e Backend}

A comunicação entre o frontend e o backend será realizada por meio de uma API REST, utilizando o protocolo HTTP.

O fluxo básico de comunicação será:

\begin{center}
\textbf{Usuário}\\
$\downarrow$\\
\textbf{Frontend}\\
$\downarrow$\\
\textbf{Requisição HTTP}\\
$\downarrow$\\
\textbf{API REST (Backend)}\\
$\downarrow$\\
\textbf{Controller}\\
$\downarrow$\\
\textbf{Service}\\
$\downarrow$\\
\textbf{Repository}\\
$\downarrow$\\
\textbf{Banco de Dados}\\
$\uparrow$\\
\textbf{Repository}\\
$\uparrow$\\
\textbf{Service}\\
$\uparrow$\\
\textbf{Controller}\\
$\uparrow$\\
\textbf{Resposta HTTP (JSON)}\\
$\uparrow$\\
\textbf{Frontend}\\
$\uparrow$\\
\textbf{Usuário}
\end{center}

\section{Organização em camadas}

O backend do Acervo Manager será organizado em camadas, com o objetivo de separar as responsabilidades de cada parte do sistema e facilitar a manutenção, a compreensão e a evolução do software.

A arquitetura será composta principalmente pelas camadas \textbf{Controller}, \textbf{Service} e \textbf{Repository}, que trabalharão de forma integrada para processar as requisições e realizar as operações no banco de dados.

\subsection{Controller}

Será responsável por receber as requisições HTTP enviadas pelo frontend, interpretar os dados recebidos e encaminhar as solicitações para a camada Service. Também será responsável por retornar as respostas HTTP ao cliente.

\subsection{Service}

Concentrará as regras de negócio do sistema, como verificar a disponibilidade de um livro antes de realizar um empréstimo, validar as condições de uma devolução e coordenar as operações necessárias.

\subsection{Repository}

Será responsável pela comunicação com o banco de dados MySQL, realizando operações como cadastro, consulta, atualização e exclusão de informações.

Essa separação permite que cada camada tenha uma função específica, reduzindo o acoplamento entre os componentes e facilitando a realização de alterações e manutenções no sistema.

O fluxo de comunicação entre as camadas será:

\begin{center}
\textbf{Frontend}
$\rightarrow$
\textbf{Controller}
$\rightarrow$
\textbf{Service}
$\rightarrow$
\textbf{Repository}
$\rightarrow$
\textbf{Banco de Dados}
\end{center}

Após a operação ser concluída, o resultado seguirá o caminho inverso até chegar ao frontend.

\section{Tratamento de erros e validações}

O Acervo Manager contará com mecanismos de validação e tratamento de erros para garantir maior confiabilidade nas operações realizadas pelo sistema. As validações terão como objetivo impedir o cadastro de informações inválidas e verificar se as operações solicitadas podem ser executadas corretamente.

No frontend, serão realizadas validações básicas nos formulários, como a verificação de campos obrigatórios, formatos de dados e preenchimento correto das informações. No backend, os dados recebidos também serão validados, garantindo que as regras de negócio sejam respeitadas independentemente da interface utilizada.

Entre as principais validações e situações de erro que poderão ser tratadas, destacam-se:

\begin{itemize}
    \item Impedir o cadastro de livros ou usuários com campos obrigatórios vazios;
    \item Verificar se o livro ou usuário informado existe no sistema;
    \item Impedir o empréstimo de um livro que não esteja disponível;
    \item Verificar se um empréstimo existe antes de realizar uma devolução;
    \item Impedir operações com dados inválidos ou incompatíveis;
    \item Retornar mensagens informativas ao usuário quando uma operação não puder ser concluída.
\end{itemize}

Dessa forma, o tratamento de erros e as validações contribuirão para a integridade dos dados, a segurança das operações e uma melhor experiência de uso do sistema.

\end{document}
